%==================================================================
% Ini adalah kata pengantar
%==================================================================

%% DILARANG EDIT BAGIAN INI
\clearpage
\phantomsection
\addcontentsline{toc}{chapter}{KATA PENGANTAR}
\begin{center}
    \textbf{\large KATA PENGANTAR}\\[3em]
\end{center}
%% DILARANG EDIT BAGIAN INI

%% Edit bagian ini sesuai kebutuhan
Puji syukur kehadirat Allah Swt. atas segala rahmat dan karunia-Nya, sehingga Tugas Akhir Skripsi ini dapat diselesaikan. Tugas Akhir ini dapat diselesaikan tidak lepas dari peran dan dukungan berbagai pihak. Berkenaan dengan hal tersebut, penulis menyampaikan ucapan terima kasih kepada:

\begin{enumerate}
    \item {\dekan} selaku Dekan Fakultas Teknik Universitas Negeri Yogyakarta yang telah memberikan dukungan dan fasilitas dalam proses penyelesaian studi.
    \item {\koorprodi} selaku Ketua Program Studi Teknologi Informasi sekaligus Dosen Pembimbing Akademik yang telah memberikan bimbingan, arahan, dan motivasi selama proses perkuliahan.
    \item {\pembimbing} selaku dosen pembimbing yang telah banyak memberikan semangat, dorongan, bimbingan, dan masukan selama penyusunan Tugas Akhir Skripsi ini.
    \item Muhammad Resa Arif Yudianto, M.Kom. selaku validator instrumen penelitian yang telah memberikan bantuan, masukan, dan saran yang sangat berharga bagi penulis.
    \item Ir. Ardy Seto Priambodo, S.T., M.Eng., yang telah menyediakan format Tugas Akhir LaTeX yang memudahkan penulis dalam menyusun Tugas Akhir Skripsi ini.
    \item Seluruh dosen dan staf pengajar di Fakultas Teknik Universitas Negeri Yogyakarta yang telah memberikan ilmu pengetahuan, wawasan, dan pengalaman berharga selama masa perkuliahan.
    \item Teman-teman anggota dan pengurus INFINITE UKM Rekayasa Teknologi yang telah memberikan semangat dan dukungan selama proses pengembangan dan penyusunan proyek Tugas Akhir Skripsi ini.
    \item Semua pihak, secara langsung maupun tidak langsung, yang tidak dapat disebutkan di sini atas bantuan dan perhatiannya selama penyusunan Tugas Akhir Skripsi ini.
\end{enumerate}

Penulis menyadari bahwa Tugas Akhir Skripsi ini masih memiliki banyak potensi untuk dikembangkan dan disempurnakan. Oleh karena itu, penulis mengharapkan kritik dan saran yang membangun dari berbagai pihak guna kelanjutan pengembangan hasil penelitian di masa yang akan datang. Penulis juga berharap Tugas Akhir Skripsi ini dapat memberikan manfaat dan dorongan bagi pengembangan ilmu pengetahuan, khususnya di bidang Teknologi Informasi modern.
%% Edit bagian ini sesuai kebutuhan

%% DILARANG EDIT BAGIAN INI
\begin{flushright}
    Yogyakarta, \tglpengesahan\\[1.25cm]
    \penulis \\
    \nim
\end{flushright}
%% DILARANG EDIT BAGIAN INI